% Section 8 — Conclusion
\section{Conclusion}
\label{sec:conclusion}

\subsection{Answer to the research question}
\label{sec:answer}
Effectiveness does not track sophistication monotonically. Moving from
traditional methods to neural forecasting improves forecast accuracy
consistently (the LSTM records the lowest error on both datasets) but
improves inventory outcomes only conditionally: substantially on sparse
intermittent demand, and not at all on dense demand under the default cost
structure. The model ladder purchases capability throughout and value in
specific settings.

\subsection{Main evidence}
\label{sec:main-evidence}
M5: LSTM lowest forecast error (MASE 1.316) and lowest inventory cost
(152.83), ranking first in 25 of 27 policies. Store: LSTM lowest forecast
error (MASE 0.978) but fourth-lowest inventory cost (2,247.46), behind Moving
Average (2,084.50), SES, and Seasonal Naive; policy wins split 9/7/6/5.
Ninety-one paired tests separate formal significance from magnitude (M5
LSTM-TSB: $p\approx2\times 10^{-10}$ at $d_z=-0.031$; Store LSTM-MA:
$d_z=-0.29$), subject to the dependence qualification in
\cref{sec:stats-method}. H1 is confirmed, H2 confirmed with qualification, H3
rejected in monotone form, and H4 rejected in strong form.

\subsection{Practical takeaway}
\label{sec:takeaway}
For sparse, intermittent assortments, the learner is supported: the forecast
advantage reaches cost and persists across policies. For dense, smooth
assortments, a moving average should be tested against costlier options
under the applicable cost ratio before investment, and any winner should be
re-tested when lead times or penalties change. In all cases, procurement
should require inventory-level evaluation rather than accepting an accuracy
leaderboard.

\subsection{Final research conclusion}
\label{sec:final-word}
Forecasting skill and decision value are related but distinct quantities,
joined by demand structure and cost asymmetry. This study measured both,
located where they correspond and where they part, stress-tested the results
across 27 policies, and froze every number for audit. Whether the
lowest-error forecaster is also the lowest-cost inventory input depends on
the setting; the evidence now records the settings on each side.
